\chapter{DESCRIPCIÓN DE LOS HECHOS}

El sistema desarrollado por Amazon utilizaba técnicas de procesamiento de lenguaje natural y modelos de aprendizaje automático supervisado, entrenados con los currículums recibidos por la empresa durante un periodo de aproximadamente diez años \cite{dastin2018}. Dado que la mayoría de los solicitantes de empleo en puestos técnicos durante ese periodo eran hombres —reflejo de la conocida brecha de género en el sector tecnológico—, el algoritmo aprendió a asociar patrones lingüísticos y curriculares ``masculinos'' con el éxito en la contratación \cite{chen2023}.

En la práctica, esto se tradujo en que el sistema penalizaba los currículums que incluían la palabra ``women's'' (por ejemplo, en expresiones como ``capitana del club de ajedrez femenino''), y rebajaba la puntuación de las candidatas graduadas de universidades exclusivamente femeninas, aunque su formación técnica fuera igual o superior a la de otros candidatos \cite{dastin2018,infobae2018}. El sistema no discriminaba explícitamente por sexo —esta variable no era un input directo del modelo—, sino que aprendió a utilizar variables proxy, es decir, características indirectamente correlacionadas con el género, para replicar patrones históricos de contratación \cite{raghavan2020}.

Los datos involucrados eran, en esencia, información curricular: formación académica, experiencia laboral, terminología empleada en la redacción del currículum y afiliaciones extracurriculares. Ninguno de estos datos pertenece, en sentido estricto, a las categorías especiales de datos personales reguladas por el RGPD (como el origen étnico o la salud), pero su tratamiento conjunto permitió inferir de forma indirecta una característica protegida —el género— y traducirla en un trato desigual, lo que constituye un caso claro de discriminación indirecta o ``proxy discrimination'' \cite{kochling2020}.

Amazon intentó corregir el sesgo neutralizando la penalización explícita de términos como ``women's'', pero, según las fuentes consultadas por Reuters, la compañía no pudo garantizar que el sistema no desarrollara otras formas más sutiles de discriminación, motivo por el cual decidió abandonar definitivamente el proyecto en 2017 \cite{dastin2018}.
